\documentclass{beamer}
\usetheme{Madrid}
\usecolortheme{default}
\usepackage[utf8]{inputenc}
\usepackage{amsmath}
\usepackage{amsfonts}
\usepackage{amssymb}
\usepackage{tikz}
\usepackage{xcolor}
\usepackage{hyperref}

\title{CRISC Exam - Domain 4: Technology \& Security}
\subtitle{Technology Risk Questions 43--82 with Explanations}
\author{Compiled from ISACA CRISC Review Manual 8th Edition}
\date{\today}

\begin{document}
	
	\begin{frame}
		\titlepage
	\end{frame}
	
	\begin{frame}{Outline}
		\tableofcontents
	\end{frame}
	
	% ================ SECTION 3: SYSTEM DEVELOPMENT LIFE CYCLE ================
	
	\section{System Development Life Cycle}
	
	\begin{frame}{Q43: SDLC Purpose}
		\begin{block}{Question}
			Why should effective management of systemic risk be accomplished through a rigorous, project-based approach to system development (SDLC)?
		\end{block}
		\begin{enumerate}[(A)]
			\item To eliminate the need for risk assessment
			\item To ensure deliverables meet business needs and reduce systemic risk
			\item To guarantee project success
			\item To replace change management
		\end{enumerate}
	\end{frame}
	
	\begin{frame}{Q43: Explanation}
		\begin{block}{Correct Answer: B}
			Complex systems can deliver business value but also can be sources of significant risk. Effective management of systemic risk is best accomplished through a rigorous and comprehensive project-based approach to system development, commonly called an SDLC.
		\end{block}
		\begin{block}{Option Analysis}
			\begin{itemize}
				\item \textbf{A - Incorrect:} SDLC complements risk assessment
				\item \textbf{B - Correct:} SDLC reduces systemic risk
				\item \textbf{C - Incorrect:} No methodology guarantees success
				\item \textbf{D - Incorrect:} SDLC works with change management
			\end{itemize}
		\end{block}
	\end{frame}
	
	\begin{frame}{Q44: SDLC - Initiation Phase}
		\begin{block}{Question}
			During which SDLC phase is the need for an IT system expressed and endorsed by a sponsor, and the purpose and scope documented?
		\end{block}
		\begin{enumerate}[(A)]
			\item Initiation
			\item Development or acquisition
			\item Implementation
			\item Operation or maintenance
		\end{enumerate}
	\end{frame}
	
	\begin{frame}{Q44: Explanation}
		\begin{block}{Correct Answer: A}
			In Phase 1 (Initiation), the need for an IT system is expressed and endorsed by a sponsor, and the purpose and scope of the anticipated IT system is documented. Risk identified within business processes is used to develop system requirements, including security requirements and a security concept of operations.
		\end{block}
		\begin{block}{Option Analysis}
			\begin{itemize}
				\item \textbf{A - Correct:} Initiation documents purpose and scope
				\item \textbf{B - Incorrect:} Development involves design and construction
				\item \textbf{C - Incorrect:} Implementation configures and tests security features
				\item \textbf{D - Incorrect:} Operation is when the system performs its functions
			\end{itemize}
		\end{block}
	\end{frame}
	
	\begin{frame}{Q45: Security in SDLC}
		\begin{block}{Question}
			Why should security requirements be addressed early in the SDLC?
		\end{block}
		\begin{enumerate}[(A)]
			\item Security is less important later in the project
			\item It is considerably easier and less expensive to build security in at the onset than to add it later
			\item Security requirements are only needed for external-facing systems
			\item Security can always be added without additional cost
		\end{enumerate}
	\end{frame}
	
	\begin{frame}{Q45: Explanation}
		\begin{block}{Correct Answer: B}
			Security is a necessity in all modern business processes, and it is considerably easier and less expensive to address security requirements by building them into the project timeline at the onset than attempting to add them later. Security requirements should be specified, validated, and tested in the same manner as other requirements.
		\end{block}
		\begin{block}{Option Analysis}
			\begin{itemize}
				\item \textbf{A - Incorrect:} Security remains critical throughout
				\item \textbf{B - Correct:} Building security in early is cheaper and easier
				\item \textbf{C - Incorrect:} All systems require security
				\item \textbf{D - Incorrect:} Adding security later is more expensive
			\end{itemize}
		\end{block}
	\end{frame}
	
	\begin{frame}{Q46: Key Security Tasks in SDLC}
		\begin{block}{Question}
			Which of the following is a key security task to perform during the SDLC to document the risk associated with the development of a new program?
		\end{block}
		\begin{enumerate}[(A)]
			\item Only reviewing the final code
			\item Security categorization of the proposed system and privacy impact assessment
			\item Ignoring vendor security assessments
			\item Skipping secure coding practices
		\end{enumerate}
	\end{frame}
	
	\begin{frame}{Q46: Explanation}
		\begin{block}{Correct Answer: B}
			Key security tasks during the SDLC include security categorization of the proposed system (CIA requirements and business impact), privacy impact assessment (sensitive data processed, applicable laws), vendor/supplier security or risk assessment, security training for development staff, and awareness of vulnerabilities in the selected technology.
		\end{block}
		\begin{block}{Option Analysis}
			\begin{itemize}
				\item \textbf{A - Incorrect:} Code review is one of many activities
				\item \textbf{B - Correct:} Security categorization and PIA are key tasks
				\item \textbf{C - Incorrect:} Vendor assessments are important
				\item \textbf{D - Incorrect:} Secure coding is mandated and enforced
			\end{itemize}
		\end{block}
	\end{frame}
	
	\begin{frame}{Q47: Project Risk vs. System Risk}
		\begin{block}{Question}
			What is the distinction between project risk and system risk?
		\end{block}
		\begin{enumerate}[(A)]
			\item There is no distinction; they are the same
			\item Project risk deals with failing to deliver on time and within budget, while system risk relates to deliverables not meeting organizational needs
			\item Project risk only applies to IT projects
			\item System risk is only relevant after implementation
		\end{enumerate}
	\end{frame}
	
	\begin{frame}{Q47: Explanation}
		\begin{block}{Correct Answer: B}
			The risk practitioner should be aware of the distinction between project risk, which deals with the consequences of failing to deliver the promised results on time and within budget, and the risk that the deliverables may not meet the needs of the organization. Projects can deliver exactly what was requested yet fail because the original request was misguided.
		\end{block}
		\begin{block}{Option Analysis}
			\begin{itemize}
				\item \textbf{A - Incorrect:} There is a clear distinction
				\item \textbf{B - Correct:} Project risk vs. system risk
				\item \textbf{C - Incorrect:} Project risk applies to all projects
				\item \textbf{D - Incorrect:} System risk applies throughout the SDLC
			\end{itemize}
		\end{block}
	\end{frame}
	
	\begin{frame}{Q48: Waterfall vs. Agile}
		\begin{block}{Question}
			Which methodology is almost always better at delivering functionality of value to the user community?
		\end{block}
		\begin{enumerate}[(A)]
			\item Waterfall
			\item Agile
			\item Both are equally effective
			\item Neither is effective
		\end{enumerate}
	\end{frame}
	
	\begin{frame}{Q48: Explanation}
		\begin{block}{Correct Answer: B}
			Because of their flexibility toward requirements, Agile methodologies are almost always better at delivering functionality of value to the user community. However, this flexibility makes controlling cost more complicated than under a Waterfall approach, where a precise understanding of what will be needed makes it possible to predict associated costs.
		\end{block}
		\begin{block}{Option Analysis}
			\begin{itemize}
				\item \textbf{A - Incorrect:} Waterfall is better for cost prediction
				\item \textbf{B - Correct:} Agile is better for delivering business value
				\item \textbf{C - Incorrect:} Each has distinct strengths
				\item \textbf{D - Incorrect:} Both are effective in different contexts
			\end{itemize}
		\end{block}
	\end{frame}
	
	\begin{frame}{Q49: Changeover Techniques - Parallel}
		\begin{block}{Question}
			Which changeover technique operates both the new and old systems simultaneously, providing the safest means of changeover if rollback is needed?
		\end{block}
		\begin{enumerate}[(A)]
			\item Parallel
			\item Phased
			\item Abrupt
			\item Direct
		\end{enumerate}
	\end{frame}
	
	\begin{frame}{Q49: Explanation}
		\begin{block}{Correct Answer: A}
			Parallel changeover provides the safest and quickest means of changeover in case the changes need to be rolled back. Since the current system is fully operational, the new system can be stopped at any point where results are not as intended. Disadvantages include the cost of maintaining and monitoring both systems and ensuring data consistency.
		\end{block}
		\begin{block}{Option Analysis}
			\begin{itemize}
				\item \textbf{A - Correct:} Parallel runs both systems simultaneously
				\item \textbf{B - Incorrect:} Phased replaces individual components
				\item \textbf{C - Incorrect:} Abrupt is a single-instant movement
				\item \textbf{D - Incorrect:} Direct is not a standard changeover term
			\end{itemize}
		\end{block}
	\end{frame}
	
	\begin{frame}{Q50: Changeover Techniques - Abrupt}
		\begin{block}{Question}
			Which changeover technique is the riskiest due to the potential for lost business processing if reversion is needed?
		\end{block}
		\begin{enumerate}[(A)]
			\item Parallel
			\item Phased
			\item Abrupt
			\item Incremental
		\end{enumerate}
	\end{frame}
	
	\begin{frame}{Q50: Explanation}
		\begin{block}{Correct Answer: C}
			Abrupt changeover is a single-instant movement from the old system to the new system, with the old system immediately taken offline. It is the riskiest of the changeover processes because of the potential for lost opportunities in business processing if it becomes necessary to revert to the earlier system. It is most prevalent when rollback is relatively assured or the impact of lost processing is minor.
		\end{block}
		\begin{block}{Option Analysis}
			\begin{itemize}
				\item \textbf{A - Incorrect:} Parallel is the safest
				\item \textbf{B - Incorrect:} Phased reduces risk by gradual rollout
				\item \textbf{C - Correct:} Abrupt is the riskiest
				\item \textbf{D - Incorrect:} Incremental is a form of phased
			\end{itemize}
		\end{block}
	\end{frame}
	
	\begin{frame}{Q51: Data Migration Considerations}
		\begin{block}{Question}
			Which of the following is a key consideration when migrating data from one system to another?
		\end{block}
		\begin{enumerate}[(A)]
			\item Data completeness, integrity, and consistency
			\item Only the speed of migration
			\item The color of the user interface
			\item The number of project team members
		\end{enumerate}
	\end{frame}
	
	\begin{frame}{Q51: Explanation}
		\begin{block}{Correct Answer: A}
			Key data conversion considerations include completeness of data (verifying total records), data integrity (not altered except as required by conversion), storage and security of data under conversion, data consistency, and business continuity. The risk practitioner should assess the migration process to identify likely areas of concern.
		\end{block}
		\begin{block}{Option Analysis}
			\begin{itemize}
				\item \textbf{A - Correct:} Completeness, integrity, consistency
				\item \textbf{B - Incorrect:} Speed is secondary to data quality
				\item \textbf{C - Incorrect:} Interface color is irrelevant
				\item \textbf{D - Incorrect:} Team size is not a data consideration
			\end{itemize}
		\end{block}
	\end{frame}
	
	\begin{frame}{Q52: Post-Implementation Review}
		\begin{block}{Question}
			Why is a timely post-implementation review important after a project?
		\end{block}
		\begin{enumerate}[(A)]
			\item To assign blame for project failures
			\item To capture lessons learned that can be applied to future projects
			\item To immediately start the next project
			\item To reduce the project budget
		\end{enumerate}
	\end{frame}
	
	\begin{frame}{Q52: Explanation}
		\begin{block}{Correct Answer: B}
			Whether an implementation project succeeds or fails, a timely post-implementation review offers the best opportunity to capture lessons learned so they can be applied to future projects. Questions include what went well, whether risk was brought within acceptable levels, whether requirements were met, and whether methodologies were followed.
		\end{block}
		\begin{block}{Option Analysis}
			\begin{itemize}
				\item \textbf{A - Incorrect:} Review is for learning, not blame
				\item \textbf{B - Correct:} Lessons learned for future projects
				\item \textbf{C - Incorrect:} Review is about reflection, not rushing ahead
				\item \textbf{D - Incorrect:} Review does not reduce budget
			\end{itemize}
		\end{block}
	\end{frame}
	
	% ================ SECTION 4: DATA LIFE CYCLE & PORTFOLIO MANAGEMENT ================
	
	\section{Data Life Cycle \& Portfolio Management}
	
	\begin{frame}{Q53: Data Life Cycle Phases}
		\begin{block}{Question}
			Which of the following correctly lists the phases of the data management life cycle?
		\end{block}
		\begin{enumerate}[(A)]
			\item Create, store, use, share, archive, destroy
			\item Identify, assess, respond, monitor
			\item Plan, build, run, monitor
			\item Budget, staff, train, audit
		\end{enumerate}
	\end{frame}
	
	\begin{frame}{Q53: Explanation}
		\begin{block}{Correct Answer: A}
			Data is commonly viewed as following a six-phase life cycle: create, store, use, share, archive, and destroy. In practice, this life cycle contains many caveats—creation may refer to original collection or synthesis from other data, and archiving or destruction may be mandated or prohibited by law or contract.
		\end{block}
		\begin{block}{Option Analysis}
			\begin{itemize}
				\item \textbf{A - Correct:} Create, store, use, share, archive, destroy
				\item \textbf{B - Incorrect:} These are risk management phases
				\item \textbf{C - Incorrect:} These are IT management phases
				\item \textbf{D - Incorrect:} These are management activities
			\end{itemize}
		\end{block}
	\end{frame}
	
	\begin{frame}{Q54: Data Ownership Risk}
		\begin{block}{Question}
			Why is a lack of clear data ownership a significant risk in the data life cycle?
		\end{block}
		\begin{enumerate}[(A)]
			\item It makes it harder to verify what requirements apply and obtain approval for archiving, sharing, or destruction
			\item It improves data security
			\item It reduces storage costs
			\item It speeds up data processing
		\end{enumerate}
	\end{frame}
	
	\begin{frame}{Q54: Explanation}
		\begin{block}{Correct Answer: A}
			A lack of clear ownership of data makes it harder to verify what requirements apply to the data and obtain approval to take steps such as archiving, sharing, or destruction. The risk practitioner should review and assess the data ownership and management process, including protection from improper disclosure, modification, or deletion.
		\end{block}
		\begin{block}{Option Analysis}
			\begin{itemize}
				\item \textbf{A - Correct:} Unclear ownership hinders verification and approval
				\item \textbf{B - Incorrect:} Lack of ownership weakens security
				\item \textbf{C - Incorrect:} Ownership is unrelated to cost
				\item \textbf{D - Incorrect:} Ownership does not affect processing speed
			\end{itemize}
		\end{block}
	\end{frame}
	
	\begin{frame}{Q55: Data Validation - Allowlist vs. Denylist}
		\begin{block}{Question}
			Which data validation approach is preferred in environments where input data validation is based on static or infrequently changing values?
		\end{block}
		\begin{enumerate}[(A)]
			\item Denylist
			\item Allowlist
			\item Random sampling
			\item Manual review
		\end{enumerate}
	\end{frame}
	
	\begin{frame}{Q55: Explanation}
		\begin{block}{Correct Answer: B}
			An allowlist is the preferred approach in environments where input data validation is based on static or infrequently changing values, where it is easier to anticipate what values would be reasonable than to anticipate what values might be malicious. A common library for allowlisting can ensure consistent data validation across multiple applications.
		\end{block}
		\begin{block}{Option Analysis}
			\begin{itemize}
				\item \textbf{A - Incorrect:} Denylists are useful when the range of valid values is extremely broad
				\item \textbf{B - Correct:} Allowlists are preferred for static values
				\item \textbf{C - Incorrect:} Random sampling is not a validation approach
				\item \textbf{D - Incorrect:} Manual review is not automated validation
			\end{itemize}
		\end{block}
	\end{frame}
	
	\begin{frame}{Q56: Data Loss Prevention - False Positive}
		\begin{block}{Question}
			What occurs when a data loss prevention (DLP) system incorrectly flags nonsensitive data as a potential data leak?
		\end{block}
		\begin{enumerate}[(A)]
			\item True positive
			\item False positive
			\item False negative
			\item True negative
		\end{enumerate}
	\end{frame}
	
	\begin{frame}{Q56: Explanation}
		\begin{block}{Correct Answer: B}
			A false positive in DLP occurs when nonsensitive data is incorrectly flagged as a potential data leak. This can lead to disrupted workflows, wasted time for cybersecurity teams, and potentially lowered confidence in the DLP system. A false negative occurs when sensitive data is mistakenly allowed to leave the organization.
		\end{block}
		\begin{block}{Option Analysis}
			\begin{itemize}
				\item \textbf{A - Incorrect:} A true positive correctly identifies sensitive data
				\item \textbf{B - Correct:} Nonsensitive data incorrectly flagged
				\item \textbf{C - Incorrect:} A false negative misses sensitive data
				\item \textbf{D - Incorrect:} A true negative correctly allows nonsensitive data
			\end{itemize}
		\end{block}
	\end{frame}
	
	\begin{frame}{Q57: Data Loss Prevention - False Negative}
		\begin{block}{Question}
			What is the primary risk of a false negative in a data loss prevention (DLP) system?
		\end{block}
		\begin{enumerate}[(A)]
			\item Disrupted workflows
			\item Sensitive data being mistakenly allowed to leave the organization
			\item Wasted time for cybersecurity teams
			\item Lowered confidence in the DLP system
		\end{enumerate}
	\end{frame}
	
	\begin{frame}{Q57: Explanation}
		\begin{block}{Correct Answer: B}
			A false negative in DLP occurs when sensitive data is mistakenly allowed to leave an organization, increasing the risk of data breaches or compliance violations. Disrupted workflows, wasted time, and lowered confidence are consequences of false positives, not false negatives.
		\end{block}
		\begin{block}{Option Analysis}
			\begin{itemize}
				\item \textbf{A - Incorrect:} This is a consequence of false positives
				\item \textbf{B - Correct:} Sensitive data leaves the organization
				\item \textbf{C - Incorrect:} This is a consequence of false positives
				\item \textbf{D - Incorrect:} This is a consequence of false positives
			\end{itemize}
		\end{block}
	\end{frame}
	
	\begin{frame}{Q58: Project vs. Program}
		\begin{block}{Question}
			What distinguishes a project from a program?
		\end{block}
		\begin{enumerate}[(A)]
			\item A project has a defined end-state, while a program is an ongoing activity
			\item A program has a defined end-state, while a project is ongoing
			\item Projects are always larger than programs
			\item There is no distinction
		\end{enumerate}
	\end{frame}
	
	\begin{frame}{Q58: Explanation}
		\begin{block}{Correct Answer: A}
			A project is a series of tasks that work toward a common purpose with a defined end-state. Activities that continue indefinitely are considered programs, which may encompass one or more projects in the course of pursuing their recurring or ongoing goals.
		\end{block}
		\begin{block}{Option Analysis}
			\begin{itemize}
				\item \textbf{A - Correct:} Project has an end; program is ongoing
				\item \textbf{B - Incorrect:} The reverse is true
				\item \textbf{C - Incorrect:} Size is not the distinction
				\item \textbf{D - Incorrect:} There is a clear distinction
			\end{itemize}
		\end{block}
	\end{frame}
	
	\begin{frame}{Q59: Project Failure Reasons}
		\begin{block}{Question}
			Which of the following is a common reason IT projects fail?
		\end{block}
		\begin{enumerate}[(A)]
			\item Clear and stable requirements
			\item Unclear or changing requirements and scope creep
			\item Adequate resourcing
			\item Effective communication
		\end{enumerate}
	\end{frame}
	
	\begin{frame}{Q59: Explanation}
		\begin{block}{Correct Answer: B}
			IT projects fail for various reasons, including unclear or changing requirements, addition of requirements without adequate resourcing (scope creep), inadequate access to resources, problems with technology, delays in delivery, unrealistic timelines, and lack of progress reporting or effective communication.
		\end{block}
		\begin{block}{Option Analysis}
			\begin{itemize}
				\item \textbf{A - Incorrect:} Clear requirements contribute to success
				\item \textbf{B - Correct:} Unclear requirements and scope creep cause failure
				\item \textbf{C - Incorrect:} Adequate resourcing contributes to success
				\item \textbf{D - Incorrect:} Effective communication contributes to success
			\end{itemize}
		\end{block}
	\end{frame}
	
	\begin{frame}{Q60: Project Failure - Time, Quality, Cost Balance}
		\begin{block}{Question}
			When a project team is struggling to remain on time and within budget, what is a common, dangerous shortcut?
		\end{block}
		\begin{enumerate}[(A)]
			\item Requesting additional resources
			\item Skipping critical steps such as testing
			\item Extending the timeline
			\item Reducing project scope
		\end{enumerate}
	\end{frame}
	
	\begin{frame}{Q60: Explanation}
		\begin{block}{Correct Answer: B}
			Time, quality, and cost exist in a balance. A project team struggling to remain on time and within budget while maintaining access to necessary resources may skip critical steps such as testing to attempt to balance these priorities. Such actions create a temporary illusion of success but may lead to production failures, inaccurate reporting, or missing functionality.
		\end{block}
		\begin{block}{Option Analysis}
			\begin{itemize}
				\item \textbf{A - Incorrect:} Requesting resources is a positive approach
				\item \textbf{B - Correct:} Skipping testing is a dangerous shortcut
				\item \textbf{C - Incorrect:} Extending timeline is a legitimate option
				\item \textbf{D - Incorrect:} Reducing scope is a legitimate option
			\end{itemize}
		\end{block}
	\end{frame}
	
	\begin{frame}{Q61: Portfolio Management Goal}
		\begin{block}{Question}
			What is the primary goal of portfolio management?
		\end{block}
		\begin{enumerate}[(A)]
			\item To eliminate all projects
			\item To optimize business value and create greater oversight under one accountable authority
			\item To increase the number of projects
			\item To reduce the IT budget
		\end{enumerate}
	\end{frame}
	
	\begin{frame}{Q61: Explanation}
		\begin{block}{Correct Answer: B}
			The goal of portfolio management is to optimize business value and create a greater degree of oversight under one accountable authority who is also appropriately empowered to solve problems as they arise. A single portfolio may include projects using different methodologies.
		\end{block}
		\begin{block}{Option Analysis}
			\begin{itemize}
				\item \textbf{A - Incorrect:} Portfolio management does not eliminate projects
				\item \textbf{B - Correct:} Optimize value and oversight
				\item \textbf{C - Incorrect:} Quantity is not the goal
				\item \textbf{D - Incorrect:} Budget reduction is not the primary goal
			\end{itemize}
		\end{block}
	\end{frame}
	
	% ================ SECTION 5: TECHNOLOGY RESILIENCE & DISASTER RECOVERY ================
	
	\section{Technology Resilience \& Disaster Recovery}
	
	\begin{frame}{Q62: Cloud and Disaster Recovery}
		\begin{block}{Question}
			How has the shift to cloud computing changed traditional disaster recovery responsibilities for customer enterprises?
		\end{block}
		\begin{enumerate}[(A)]
			\item Enterprises now have more control over recovery
			\item Enterprises can often do nothing to recover until their cloud providers regain functionality
			\item Disaster recovery is no longer needed
			\item Enterprises now own all recovery infrastructure
		\end{enumerate}
	\end{frame}
	
	\begin{frame}{Q62: Explanation}
		\begin{block}{Correct Answer: B}
			Generally lacking hardware and other physical infrastructure of their own, cloud-reliant enterprises can generally do nothing to recover their operations until their hosts regain a level of operating capability sufficient to meet their expectations. Enterprises using software as a service will be unable to resume production until the software returns to normal operations.
		\end{block}
		\begin{block}{Option Analysis}
			\begin{itemize}
				\item \textbf{A - Incorrect:} Cloud generally reduces recovery control
				\item \textbf{B - Correct:} Enterprises wait for provider recovery
				\item \textbf{C - Incorrect:} Disaster recovery is still needed
				\item \textbf{D - Incorrect:} Cloud providers own the infrastructure
			\end{itemize}
		\end{block}
	\end{frame}
	
	\begin{frame}{Q63: Most Important Role in Technology Resilience}
		\begin{block}{Question}
			What is the most important role the risk practitioner can plan as it relates to technology resilience?
		\end{block}
		\begin{enumerate}[(A)]
			\item Personally recover all systems
			\item Ensure the risk assumed through cloud technologies is well documented, fully understood, and formally accepted by management
			\item Eliminate all cloud usage
			\item Purchase additional insurance
		\end{enumerate}
	\end{frame}
	
	\begin{frame}{Q63: Explanation}
		\begin{block}{Correct Answer: B}
			The most important role that the risk practitioner can plan as it relates to technology resilience is to ensure that the risk assumed by the enterprise through its particular combination of cloud technologies and services is well documented, fully understood, and formally accepted by management.
		\end{block}
		\begin{block}{Option Analysis}
			\begin{itemize}
				\item \textbf{A - Incorrect:} Risk practitioners do not personally recover systems
				\item \textbf{B - Correct:} Document, understand, and accept cloud risk
				\item \textbf{C - Incorrect:} Eliminating cloud is often infeasible
				\item \textbf{D - Incorrect:} Insurance is one option, not the main role
			\end{itemize}
		\end{block}
	\end{frame}
	
	\begin{frame}{Q64: SLA Misunderstanding}
		\begin{block}{Question}
			What is a common misunderstanding about service level agreements (SLAs) in cloud environments?
		\end{block}
		\begin{enumerate}[(A)]
			\item SLAs guarantee uptime
			\item SLAs ensure agreed-upon compensation for each instance in which the guaranteed level is not met, not guaranteed uptime
			\item SLAs eliminate all risk
			\item SLAs are not legally binding
		\end{enumerate}
	\end{frame}
	
	\begin{frame}{Q64: Explanation}
		\begin{block}{Correct Answer: B}
			Problems arise when managers have an imperfect understanding of how much risk the enterprise is taking on, such as mistakenly believing that an SLA ensures particular levels of uptime rather than ensuring agreed-upon compensation for each instance in which the guaranteed level is not met. A credit in the next billing cycle may not be a sufficient offset for an outage that puts the enterprise at risk of collapse.
		\end{block}
		\begin{block}{Option Analysis}
			\begin{itemize}
				\item \textbf{A - Incorrect:} SLAs provide compensation, not uptime guarantees
				\item \textbf{B - Correct:} SLAs ensure compensation, not uptime
				\item \textbf{C - Incorrect:} SLAs do not eliminate risk
				\item \textbf{D - Incorrect:} SLAs are legally binding
			\end{itemize}
		\end{block}
	\end{frame}
	
	\begin{frame}{Q65: Risk Transfer and Reputation}
		\begin{block}{Question}
			Why is risk transfer through cloud services or insurance incomplete?
		\end{block}
		\begin{enumerate}[(A)]
			\item Because contracts are always void
			\item Because the enterprise retains the impact to its reputation in the event of loss
			\item Because insurance companies never pay claims
			\item Because risk cannot be transferred at all
		\end{enumerate}
	\end{frame}
	
	\begin{frame}{Q65: Explanation}
		\begin{block}{Correct Answer: B}
			The use of cloud services is a form of risk transfer for which the enterprise retains the impact to its reputation in the event of loss. A credit in the next billing cycle may not be a sufficient offset for an outage that puts the enterprise at risk of collapse. Stakeholders rarely regard the transfer of risk as complete absolution of blame.
		\end{block}
		\begin{block}{Option Analysis}
			\begin{itemize}
				\item \textbf{A - Incorrect:} Contracts are generally valid
				\item \textbf{B - Correct:} Reputational impact remains with the enterprise
				\item \textbf{C - Incorrect:} Insurers generally pay valid claims
				\item \textbf{D - Incorrect:} Risk can be transferred, but not completely
			\end{itemize}
		\end{block}
	\end{frame}
	
	% ================ FINAL SLIDES ================
	
	\begin{frame}{SDLC \& Resilience Key Takeaways}
		\begin{block}{Essential Concepts}
			\begin{itemize}
				\item SDLC provides a structured approach to manage systemic risk
				\item Security must be built in early—it is cheaper and more effective
				\item Project risk (on time/budget) differs from system risk (meets needs)
				\item Agile delivers better business value; Waterfall enables better cost control
				\item Parallel changeover is safest; Abrupt is riskiest
				\item Data life cycle: create, store, use, share, archive, destroy
				\item Allowlists are preferred for static data validation
				\item DLP false positives disrupt workflows; false negatives allow leaks
				\item Cloud shifts recovery responsibility to providers—enterprises must document and accept the risk
			\end{itemize}
		\end{block}
	\end{frame}
	
	\begin{frame}{Final Advice}
		\begin{block}{Exam Preparation}
			\centering
			\vspace{0.3cm}
			\textbf{Remember:}
			\begin{itemize}
				\item Understand the \textit{why} behind each technology risk concept
				\item Know the difference between similar terms
				\item Understand how emerging technologies introduce new risk
				\item Practice scenario-based questions
				\item Use the ISACA mindset
			\end{itemize}
			\vspace{0.5cm}
			\textbf{Good Luck on Your CRISC Exam!}
		\end{block}
	\end{frame}
	
\end{document}